\documentclass[10pt,a4paper]{amsart}
\usepackage[margin=19mm]{geometry}
\usepackage{amsmath,amssymb,mathtools}
\usepackage[scr=rsfs,cal=euler]{mathalfa}
\usepackage{xfrac}
\usepackage[unicode,hidelinks,bookmarks=false]{hyperref}
\newcommand{\ie}{i.e.\ }
\newcommand{\cf}{cf.\ }
\newcommand{\resp}{respectively\ }
\newcommand{\Subsection}{\subsection}
\providecommand{\nobreakdash}{\nobreak\hbox{-}}
\setlength{\emergencystretch}{2em}
\multlinegap0pt
\usepackage{amssymb}
\usepackage[shortlabels]{enumitem}
\newtheorem{lemma}{Lemma}
\title{Categorical Algebra of Atomic Monoids: Presentability, Regularity, and Pretorsion Theories}
\author{Federico Campanini, Laura Cossu}
\date{Source: arXiv:2607.23144v1 (CC BY 4.0)}
\begin{document}
\maketitle

\noindent\emph{Source and licence.} Categorical Algebra of Atomic Monoids: Presentability, Regularity, and Pretorsion Theories. Version arXiv:2607.23144v1 (CC BY 4.0), distributed by arXiv under the Creative Commons Attribution 4.0 International licence, http://creativecommons.org/licenses/by/4.0/, as recorded in the arXiv licence metadata for this exact version and shown on the abstract page https://arxiv.org/abs/2607.23144v1. Full licence notice: \texttt{licenses/arxiv-2607.23144-CC-BY-4.0.txt}.\par\smallskip

\noindent\textit{Editorial scope.} Counted unit: the article's lemma lem:pullback-presentation (source0.tex lines 570--579). The article's complete original proof of it is delivered with it (source0.tex lines 581--644, one \texttt{proof} environment of 64 lines). No mathematics is written, repaired or re-derived: the statement and the proof are the article's own lines, in the article's own notation, and no retained line is substituted or re-flowed.  Retained uncounted as the source-local context the proof needs: lines 564--568.  Nothing was omitted from any retained span, so the package carries no omission record.  No retained line contains a citation, so the wrapper carries no bibliography and no bibliography entry.  No retained line contains a figure environment, an image-inclusion command or drawing code, so no image is delivered and the package declares no asset dependency.  Editorial formatting only, in addition: an AMS \texttt{amsart} preamble that restates the article's own declarations and macros used by the retained lines, namely the environments \texttt{lemma} on separate counters, together with the standard packages those lines need.  The retained environments print in this document as Lemma 1 in order of appearance, whereas the article prints the same items with the numbering of its own full document.  The complete native source of the article as distributed in the arXiv e-print for this exact version is bundled unchanged as \texttt{sources/206005/source0.tex}.
 \begin{equation}\label{eq:ABCD}
 A_i:=(a_i,a),\qquad C_i:=(c_i,c)
 \quad \text{for }i\in\{1,2,3\},\qquad 
 B:=(b,b),\qquad D:=(d,d).
 \end{equation}

\begin{lemma}\label{lem:pullback-presentation}
The pullback object \(P\) admits the presentation
$$
P\cong \left\langle A_1,A_2,A_3,B,C_1,C_2,C_3,D \ \middle|\ A_1BA_3C_2=A_2C_1C_3D \right\rangle_{\mathrm{ab}}.
$$
Under this identification,
$$
\rho(A_i)=a,\qquad \rho(B)=b,\qquad \rho(C_i)=c,\qquad \rho(D)=d.
$$
\end{lemma}

\begin{proof}
Set
$$ M:= \left\langle A_1,A_2,A_3,B,C_1,C_2,C_3,D \ \middle|\ A_1BA_3C_2=A_2C_1C_3D \right\rangle_{\mathrm{ab}}. $$ 
Let $$ \Phi:M\to P $$ be the surjective monoid homomorphism determined by the correspondences in \eqref{eq:ABCD}.
The map is well defined, since the relation
\begin{equation}\label{eq:relation-M}
A_1BA_3C_2=A_2C_1C_3D
\end{equation}
is satisfied by the corresponding elements of both $H$ and $J$. We prove that $\Phi$ is injective. We first show how equalities in $H$ and $J$ can be expressed in terms of exponent vectors. Consider two elements
$$
u=a_1^{r_1}a_2^{r_2}a_3^{r_3}b^{r_4} c_1^{r_5}c_2^{r_6}c_3^{r_7}d^{r_8} \quad \text{and} \quad v=a_1^{s_1}a_2^{s_2}a_3^{s_3}b^{s_4} c_1^{s_5}c_2^{s_6}c_3^{s_7}d^{s_8}
$$
in $H$, where all the exponents are non-negative integers, and write
$$
\mathbf r=(r_1,\ldots,r_8), \qquad \mathbf s=(s_1,\ldots,s_8).
$$
If $u=v$ in $H$, then $u$ can be transformed into $v$ by applying finitely many times, in either direction, the relations
$$
a_1b=a_2c_1, \qquad a_3c_2=c_3d.
$$
Consequently, there exist $k,\ell\in\mathbb Z$ such that
$$
\mathbf r = \mathbf s +k(-1,1,0,-1,1,0,0,0) +\ell(0,0,-1,0,0,-1,1,1).
$$
Similarly, given two elements
$$
x=a^{m_1}b^{m_2}c^{m_3}d^{m_4} \quad \text{and} \quad y=a^{n_1}b^{n_2}c^{n_3}d^{n_4}
$$ in $J$, write the vectors of their non-negative exponents as
$$
\mathbf m=(m_1,m_2,m_3,m_4), \qquad \mathbf n=(n_1,n_2,n_3,n_4).
$$
If $x=y$ in $J$, then there exists $t\in\mathbb Z$ such that
$$
\mathbf m=\mathbf n+t(-1,-1,1,1).
$$
Now let $U,V\in M$, represented by
$$
U= A_1^{r_1}A_2^{r_2}A_3^{r_3}B^{r_4} C_1^{r_5}C_2^{r_6}C_3^{r_7}D^{r_8} \quad \text{and} \quad V= A_1^{s_1}A_2^{s_2}A_3^{s_3}B^{s_4} C_1^{s_5}C_2^{s_6}C_3^{s_7}D^{s_8},
$$
where all exponents are non-negative integers, and suppose that $\Phi(U)=\Phi(V)$. Then, $\sigma(\Phi(U))=\sigma(\Phi(V))$ in $H$ and from the previous observation, we obtain
$$
(r_1,\ldots,r_8) = (s_1,\ldots,s_8) +(-k,k,-\ell,-k,k,-\ell,\ell,\ell)
$$
for suitable $k,\ell\in\mathbb Z$.

On the other hand, $\rho(\Phi(U))=\rho(\Phi(V))$ in $J$.
% The exponent vectors of these two elements of $J$ are respectively
% $$
% (r_1+r_2+r_3,\ r_4,\ r_5+r_6+r_7,\ r_8) \quad \text{and} \quad (s_1+s_2+s_3,\ s_4,\ s_5+s_6+s_7,\ s_8).
% $$
Hence there exists $t\in\mathbb Z$ such that
$$
(r_1+r_2+r_3,\ r_4,\ r_5+r_6+r_7,\ r_8) = (s_1+s_2+s_3,\ s_4,\ s_5+s_6+s_7,\ s_8) +t(-1,-1,1,1).
$$
Combining the two vector equalities, we get $k=\ell=t$. It follows that
$$
(r_1,\ldots,r_8) = (s_1,\ldots,s_8) +k(-1,1,-1,-1,1,-1,1,1).
$$
If $k \geq 0$, then there exists an element $W \in M$ such that
$$
U=W(A_2C_1C_3D)^k \quad \text{and} \quad V=W(A_1BA_3C_2)^k
$$
from which $U=V$. If $k\leq 0$ it suffices to interchange the roles of $U$ and $V$. Thus, $\Phi$ is injective.
\end{proof}

\end{document}
